% !TeX root = ../../tesis.tex
%%%%%%%%%%%%%%%%%%%%%%%%%%%%%%%%%%%%%%%%%%%%%%%%%%%%%%%%%%%%%%%%%%%
% SECCIÓN 5.5.3 — VULNERABILIDAD ESTRUCTURAL ΔE
%%%%%%%%%%%%%%%%%%%%%%%%%%%%%%%%%%%%%%%%%%%%%%%%%%%%%%%%%%%%%%%%%%%
\subsection{Vulnerabilidad Estructural (\texorpdfstring{$\Delta E$}{Delta E})}
\label{subsec:vulnerabilidad-estructural}

Los resultados de esta subsección tienen carácter preliminar. La
simulación de vulnerabilidad estructural se realizó sobre los nodos
del grafo construido por \texttt{VFTModel}, pero no fue ejecutada
sobre la capa de red integrada de \texttt{Transport-gis-mjg}; en
consecuencia, los valores reportados no han sido verificados contra
la representación completa del sistema. Se presentan como primer
acercamiento al indicador $\Delta E$, con la intención de enmarcar
el método y señalar las líneas de trabajo futuro.

A partir del valor de referencia $T = 108.92$~min
(§\ref{subsec:tiempo-viaje-promedio}), el experimento calcula la
caída relativa de eficiencia $\Delta E$ cuando se retira cada uno de
los cinco nodos de mayor $B(v)$ (§\ref{subsec:robustez-geometrica}).
La eficiencia base es $E_0 = 1/108.92$; tras remover el nodo $v$,
se recalcula el tiempo promedio $T_f$ y se obtiene
$\Delta E = (E_0 - E_f)/E_0$.

El Cuadro~\ref{tab:delta-e-top5} registra los resultados del
experimento sobre los cinco nodos seleccionados.

\begin{table}[H]
    \centering
    \caption[Impacto de remoción de nodos críticos --- $\Delta E$ exploratorio]{%
    Caída relativa de eficiencia $\Delta E$ al remover cada nodo del
    componente gigante. $E_0 = 1/108.92$; $E_f = 1/T_f$. Corrida
    exploratoria: 2026-09-02.}
    \label{tab:delta-e-top5}
    \small
    \begin{tabular}{|r|l|r|r|r|}
        \hline
        \textbf{\#} & \textbf{Nodo removido} &
        \textbf{$B(v)$} & \textbf{$T_f$ (min)} & \textbf{$\Delta E$} \\
        \hline
        1 & Tacubaya        & 0.2179 & 111.80 & 2.58\,\% \\
        2 & Mixcoac         & 0.1725 & 110.19 & 1.15\,\% \\
        3 & Hidalgo         & 0.1595 & 109.90 & 0.89\,\% \\
        4 & Balderas        & 0.1433 & 109.57 & 0.59\,\% \\
        5 & Lázaro Cárdenas & 0.1376 & 110.13 & 1.10\,\% \\
        \hline
    \end{tabular}
    \fuentefigura{Elaboración propia con \texttt{VFTModel}
    \citep{cabrera2026VFTModel}. Corrida exploratoria: 2026-09-02.}
\end{table}

Tacubaya encabeza tanto el ranking de $B(v)$ como el de $\Delta E$: su
remoción eleva el tiempo promedio de la red en 2.88~min, traducido en
una caída de eficiencia del 2.58\,\%. El dato es coherente con la
posición estructural del nodo ---cruce de las Líneas 1, 7 y 9 del
Metro--- y su papel de puente entre el corredor poniente y el centro.
Los cuatro nodos restantes producen caídas entre 0.59\,\% y 1.15\,\%,
lo que confirma que ningún nodo individual colapsa la red, pero la
pérdida acumulada de varios podría alcanzar umbrales críticos.

Los tres indicadores de Fase~3 convergen en un mismo diagnóstico
estructural. $T = 108.92$~min revela que el tiempo de viaje promedio
entre cualquier par de la \zmvm{} supera con amplitud los umbrales de
movilidad cotidiana. $B(v)$ identifica el corredor que hace posible
ese tiempo: un conjunto reducido de estaciones del Metro en el eje
centro-poniente concentra la mayor parte de los caminos mínimos de la
red. Y la simulación exploratoria de $\Delta E$ confirma que la
remoción de esos nodos, aunque no colapsa la red, la degrada de
forma mensurable ---con Tacubaya produciendo el mayor impacto
individual registrado.

El patrón es consistente con el hallazgo de Fase~1: las zonas que
registraron menor cobertura y mayor Factor de Desviación son también
las que más dependen del corredor central para cualquier trayecto de
largo aliento. La siguiente sección integra estos indicadores en un
diagnóstico territorial que localiza las zonas de déficit estructural
de la red.
